\documentclass[11pt,a4paper]{article}
\usepackage[T1]{fontenc}
\usepackage{lmodern,amsmath,amssymb,amsthm,microtype,booktabs}
\usepackage[margin=27mm]{geometry}
\usepackage[colorlinks=true,urlcolor=blue,linkcolor=blue]{hyperref}
\newtheorem{proposition}{Proposition}
\newcommand{\dd}{\mathop{}\!\mathrm d}
\title{A quantitative recoil calibration\\for finite-gravity comparisons}
\author{Alan Day \and Chengheng\thanks{AI research assistant (OpenAI).}}
\date{Working research note --- 30 October 2026}
\begin{document}
\maketitle
\begin{abstract}
We bound the finite-principal-series correction to the confluent
pure-recoil kernel of Chen, Stanford, Tang and Yang. A positive integral
gives a relative window and a nontrivial crossover. We then compare the
limiting function with the inverse transform of the reduced semiclassical
action used by Chandra, Liska and Post. Their leading actions agree along
the pure-recoil route, but their normalized prefactors differ.
These are reduced-model calculations, not a finite-gravity correlator.
The arguments are author self-reviewed drafts; no priority or independent
certification is claimed.
\end{abstract}
\section{Objects and order of limits}
Set the de Sitter radius to one and use observer proper time \(t\).
CSTY~\cite{CSTY}, Section 2.6, first takes pure recoil,
\(Gm_i,G\nu,G\nu'\to0\), with \(m,\nu'\gg1\), \(\nu'/m\ll1\).
Its fixed-average-mass mixed representation and antipodal insertions
differ from the finite vacuum packet of our two-dimensional notes.
After the small-kink approximation, equation (2.78) defines
\begin{equation}
 F_\nu(Q)={}_3F_2\left(\begin{matrix}1,1-i\nu,1+i\nu\\3/2,3/2\end{matrix};
 \frac{Q^2}{\nu^2}\right),\quad
 Q=A\cosh t,\quad A=\frac{\nu\nu'}{2m}. \label{eq:F}
\end{equation}
We take \(\nu>0\), \(0\le Q<\nu\), on the real branch below the first
branch point. Finite \(\nu\) here refers to this reduced expression;
it does not restore the errors omitted before (2.78).
Equation (2.79) takes \(\nu\to\infty\) at fixed \(Q\), giving
\begin{equation}
 H(Q)={}_1F_2(1;3/2,3/2;Q^2). \label{eq:H}
\end{equation}
Only this limit is quantified below. No finite-\(\kappa\) Virasoro
\(6j\), continuation through \(Q=\nu\), or microscopic dictionary is assumed.

\section{Positive integral and explicit error}
\begin{proposition}
For \(0<Q<\nu\), write \(r=Q/\nu\). Then
\begin{align}
 F_\nu(Q)&=\frac1{2Q}\int_0^{\pi/2}
 \frac{\sinh(2\nu\arcsin(r\cos\theta))}
 {\sqrt{1-r^2\cos^2\theta}}\,\dd\theta, \label{eq:int}\\
 H(Q)&=\frac1{2Q}\int_0^{\pi/2}\sinh(2Q\cos\theta)\,\dd\theta,\\
 1\le\frac{F_\nu(Q)}{H(Q)}
 &\le \frac{\arcsin r}{r\sqrt{1-r^2}}\,
 e^{2\nu(\arcsin r-r)}. \label{eq:bound}
\end{align}
Both functions equal one at \(Q=0\).
\end{proposition}
\begin{proof}
The beta integral with parameters \(1,1/2\) inserts
\((1)_n/(3/2)_n\) into the Gauss series, with weight
\(\dd u/(2\sqrt{1-u})\). Substitute \(u=\cos^2\theta\) and use the
elementary identity already present in CSTY (2.76):
\[
 {}_2F_1(1-i\nu,1+i\nu;3/2;z^2)
 =\frac{\sinh(2\nu\arcsin z)}{2\nu z\sqrt{1-z^2}}.
\]
Absolute convergence for \(r<1\) justifies termwise integration.
The limiting integral follows at fixed \(Q\) by dominated convergence,
or by its hyperbolic-sine series. These are standard beta-integral
operations~\cite{DLMF}.

For \(x\ge0,c\ge1\), \(\sinh(cx)\le c e^{(c-1)x}\sinh x\).
Indeed, use \(\sinh x/x=\int_0^1\cosh(xu)\dd u\) and
\(\cosh(cxu)\le e^{(c-1)x}\cosh(xu)\).
For each \(\theta\), put \(z=r\cos\theta\), \(x=2\nu z\),
\(c=\arcsin z/z\). Both \(\arcsin z/z\) and \(\arcsin z-z\)
increase on \([0,r]\). Bound the remaining factor by
\((1-r^2)^{-1/2}\). Positive integrands give the upper bound;
\(\arcsin z\ge z\) gives the lower bound.
\end{proof}
Let \(B(Q,\nu)\) be the logarithm of the right-hand side of
\eqref{eq:bound}. For \(r\to0\),
\begin{equation}
 B=\frac{Q^3}{3\nu^2}+\frac{2Q^2}{3\nu^2}
   +O\left(\frac{Q^5+Q^4}{\nu^4}\right).
\end{equation}
Thus \(0\le F_\nu/H-1\le e^B-1\) is a computable analytic bound
for these reduced functions, and
\begin{equation}
 Q=o(\nu^{2/3})\ \Longrightarrow\ F_\nu(Q)/H(Q)\to1.
 \label{eq:window}
\end{equation}
This controls relative error of a possibly exponentially large function.

\section{Crossover before the branch point}
\begin{proposition}
Fix \(r_0<1\). As \(Q\to\infty\), uniformly for \(0<r=Q/\nu\le r_0\),
\begin{align}
 F_\nu(Q)&=\frac{\sqrt\pi}{8Q^{3/2}}(1-r^2)^{-1/4}
 e^{2\nu\arcsin r}\,[1+O_{r_0}(Q^{-1})],\label{eq:saddle}\\
 H(Q)&=\frac{\sqrt\pi}{8Q^{3/2}}e^{2Q}[1+O(Q^{-1})].
\end{align}
For any fixed \(c>0\), \(Q=c\nu^{2/3}\) consequently gives
\begin{equation}
 F_\nu(Q)/H(Q)\longrightarrow e^{c^3/3}. \label{eq:cross}
\end{equation}
\end{proposition}
\begin{proof}
The positive exponential in \eqref{eq:int} is \(e^{Q f_r(\theta)}\),
where \(f_r=2\arcsin(r\cos\theta)/r\), continuously extended at \(r=0\).
Its unique maximum is at zero, with \(f_r''(0)=-2/\sqrt{1-r^2}\).
The amplitude and fixed derivatives of \(f_r\) are uniformly bounded
near zero for \(r\le r_0\). Also
\(f_r(0)-f_r(\theta)\ge2(1-\cos\theta)\), giving a uniform exponential
gap outside a fixed neighbourhood. On the scale \(\theta=O(Q^{-1/2})\),
the even Taylor expansion, Gaussian domination and that tail bound yield
relative remainder \(O_{r_0}(Q^{-1})\). The negative exponential in
the hyperbolic sine is \(O(Q^{-1})\) absolutely and exponentially smaller
than the positive saddle. The half-Gaussian factor is
\[
 \frac{1}{4Q\sqrt{1-r^2}}\,
 \frac{\sqrt\pi(1-r^2)^{1/4}}{2\sqrt Q},
\]
giving \eqref{eq:saddle}. Its \(r=0\) version gives \(H\).
At the crossover \(r=c\nu^{-1/3}\),
\(2\nu(\arcsin r-r)\to c^3/3\), and the amplitude ratio tends to one.
\end{proof}
The resulting growing-time overlap inside the reduced model is
\begin{equation}
 1\ll A\cosh t\ll\nu^{2/3}.
\end{equation}
The lower inequality permits the saddle approximation; the upper one
controls confluence at the function level. At large positive \(t\),
the upper condition is \(Ae^t/(2\nu^{2/3})\to0\), stronger than
\(Q/\nu\to0\). The latter can preserve only the leading logarithmic
action while losing relative accuracy.

The earlier heavy-background and small-kink limits require separate
control for simultaneous scalings of \(m,\nu',\nu,t\).
Even holding the leading kink relation fixed, replacing
\(\sin(\gamma/2)\) by \(\gamma/2\) changes an enhanced exponent on the
scale \(Q(\nu'/m)^2\). This is a Taylor sensitivity estimate, not a
bound on all omitted gravitational effects.

\section{A prefactor test of the action-only inverse transform}
CLP~\cite{CLP}, (3.77)--(3.83) and footnote 12, distinguish their reduced
action transform from the full finite-\(\kappa\) amplitude.
For \(A=\omega_c>0\), the former is
\begin{equation}
 I_A(t)=\int_{\mathbb R}\dd y\,
 \exp\left\{2\left[-y\,\operatorname{arsinh}(y/A)
 +\sqrt{A^2+y^2}-A+yt\right]\right\}. \label{eq:I}
\end{equation}
Its unique maximum is \(y_*=A\sinh t\).
Setting \(y=Az\), the exponent is \(A\) times a strictly concave
function, with second derivative \(-2/\cosh t\) at the saddle.
Laplace's method, uniformly on bounded real \(t\)-intervals, gives
\[
 I_A(t)=\sqrt{\pi A\cosh t}\,
 e^{2A(\cosh t-1)}[1+O(A^{-1})]\qquad(A\to\infty).
\]
The tails decay uniformly on these intervals. Thus
\begin{align}
 \frac{I_A(t)}{I_A(0)}
 &= (\cosh t)^{1/2} e^{2A(\cosh t-1)}[1+O(A^{-1})],\\
 \frac{H(A\cosh t)}{H(A)}
 &= (\cosh t)^{-3/2} e^{2A(\cosh t-1)}[1+O(A^{-1})].
\end{align}
Their ratio tends to \((\cosh t)^2\), not one.
These \(t=0\) divisions are diagnostic normalizations, not
identifications with physical disconnected correlators.

A conditional necessary test follows. Suppose a missing reduced-energy
prefactor has smooth leading form \(p(y/A)\), is nonzero at the saddle,
and preserves the dominant contour. To reproduce the normalized recoil
prefactor one would need, on the tested saddle interval,
\begin{equation}
 \frac{p(\sinh t)}{p(0)}=(\cosh t)^{-2},
 \qquad \frac{p(z)}{p(0)}=\frac1{1+z^2}. \label{eq:test}
\end{equation}
This is not a derivation or fitted prescription for a Virasoro measure.
External legs, spectral weights, determinants and other saddles may
change the premise. It gives a concrete test for a later calculation.

CLP has \(A=\pi\nu\nu'/(2\kappa\sin\theta_H)\). Along the pure-recoil
route \(\theta_H=\pi(1-m/\kappa)\), \(A\to\nu\nu'/(2m)\).
This degeneration differs from fixed nondegenerate cone angle.
No uniform finite-\(\kappa\) error in that degeneration has been
established here, so this is not a finite-\(G\) match.

\section{Numerical checks and next calculation}
The script \texttt{code/wkb\_overlap.py} uses positive integrals with
their leading exponent subtracted. Six finite-\(\nu\) integrals are
compared with the defining hypergeometric function and the explicit
bound; nine crossover points and six inverse-transform comparisons
are also evaluated. All 42 checks pass at 60-digit working precision;
crossover and inverse integrals are repeated at 85 and 80 digits.
These are empirical arithmetic checks, not interval certification.

For \(Q=\nu^{2/3}\), the computed \(\log(F_\nu/H)\) is
\[
\begin{array}{c|ccc}
\nu&10^3&10^4&10^6\\ \hline
\log(F_\nu/H)&0.337345509&0.334195668&0.333373335
\end{array}
\]
and its predicted limit is \(1/3\).
At \(t=1,A=160\), the normalized \(I/H\) ratio is \(2.380635626\);
its predicted limit is \(\cosh^2(1)=2.381097846\).
No finite list of samples proves the uniform statements.

Next derive the leading prefactor from the continued \(6j\), its vertex
factors and physical observer insertions, while following the external
energy contour. Use \eqref{eq:test} as a diagnostic, not an input.
The pure-recoil degeneration must also be controlled, or a common
finite-deficit regime chosen. The two-dimensional finite-packet theorem
does not supply that missing bridge.

\begin{thebibliography}{9}
\bibitem{CSTY} Y. Chen, D. Stanford, H. Tang and Z. Yang,
\emph{Negative shocks versus static patch holography},
\href{https://arxiv.org/html/2607.14042v1}{arXiv:2607.14042v1},
Section 2.6, (2.76)--(2.84).
\bibitem{CLP} J. Chandra, D. Liska and B. Post,
\emph{3D de Sitter from complex Virasoro TQFT: Observers and knots},
\href{https://arxiv.org/html/2610.30247v1}{arXiv:2610.30247v1},
Section 3.3, footnote 12 and Appendix C.
\bibitem{DLMF} NIST Digital Library of Mathematical Functions,
\href{https://dlmf.nist.gov/16.5}{Section 16.5},
beta-integral reductions; \href{https://dlmf.nist.gov/15.4}{Section 15.4},
elementary Gauss functions. Online formulas checked 30 October 2026.
\end{thebibliography}
\end{document}
